
\documentclass{article}
\usepackage{amsfonts}

\usepackage{eurosym}
%%%%%%%%%%%%%%%%%%%%%%%%%%%%%%%%%%%%%%%%%%%%%%%%%%%%%%%%%%%%%%%%%%%%%%%%%%%%%%%%%%%%%%%%%%%%%%%%%%%%%%%%%%%%%%%%%%%%%%%%%%%%%%%%%%%%%%%%%%%%%%%%%%%%%%%%%%%%%%%%%%%%%%%%%%%%%%%%%%%%%%%%%%%%%%%%%%%%%%%%%%%%%
\usepackage{amsmath, amssymb}

\setcounter{MaxMatrixCols}{10}
%TCIDATA{OutputFilter=Latex.dll}
%TCIDATA{Version=5.50.0.2960}
%TCIDATA{<META NAME="SaveForMode" CONTENT="1">}
%TCIDATA{BibliographyScheme=Manual}
%TCIDATA{LastRevised=Tuesday, March 11, 2025 08:12:26}
%TCIDATA{<META NAME="GraphicsSave" CONTENT="32">}

\input{tcilatex}

\begin{document}


\subsection*{1. What is a Minimal Left Ideal?}

In the context of operator algebras (or matrix algebras), a \textbf{left
ideal} $\mathcal{I} $ is a subspace of the algebra that is closed under left
multiplication by any element of the algebra. That is, if $A $ is any $2
\times 2 $ complex matrix and $B \in \mathcal{I} $, then $AB \in \mathcal{I} 
$.

A \textbf{minimal left ideal} is a left ideal that contains no nontrivial
smaller left ideals. In other words, it is the smallest possible nontrivial
subspace that is closed under left multiplication.

\begin{itemize}
\item The left ideal generated by $|\psi \rangle \langle \psi |$ consists of
all matrices of the form: 
\begin{equation*}
A\cdot |\psi \rangle \langle \psi |=%
\begin{pmatrix}
a & b \\ 
c & d%
\end{pmatrix}%
\begin{pmatrix}
1 & 0 \\ 
0 & 0%
\end{pmatrix}%
=%
\begin{pmatrix}
a & 0 \\ 
c & 0%
\end{pmatrix}%
,
\end{equation*}%
where $A=%
\begin{pmatrix}
a & b \\ 
c & d%
\end{pmatrix}%
$ is an arbitrary $2\times 2$ matrix. This subspace is 2-dimensional over $%
\mathbb{C}$ and corresponds to the space of spinors.
\end{itemize}

---

\section*{Does Such a Set Have a Name?}

The specific set of partial isometries in a \textbf{C\*-algebra} that share
the same \textbf{final projector} but have different \textbf{initial
projectors} does not have a universally standardized name in the literature.
However, it can be described in terms of well-known concepts in operator
theory and C\*-algebras. Let\^{a}\euro \U{2122}s clarify this and provide
some context.

\subsection*{1. Description of the Set}

The set of partial isometries $v $ in a C\*-algebra $\mathcal{A} $ that
share the same final projector $q = v v^* $ but have different initial
projectors $p = v^* v $ can be written as: 
\begin{equation*}
\mathcal{V}_q = \{ v \in \mathcal{A} \mid v v^* = q \}. 
\end{equation*}
This set consists of all partial isometries that map the range of their
initial projectors $p $ isometrically onto the range of the fixed final
projector $q $.

\subsection*{2. Relation to Murray-von Neumann Equivalence}

The set $\mathcal{V}_q $ is closely related to the concept of \textbf{%
Murray-von Neumann equivalence}: - Two projections $p $ and $q $ are \textbf{%
Murray-von Neumann equivalent} if there exists a partial isometry $v $ such
that $v^* v = p $ and $v v^* = q $. - The set $\mathcal{V}_q $ consists of
all partial isometries that implement this equivalence for a fixed $q $.

\subsection*{3. Possible Names and Terminology}

While there isn\^{a}\euro \U{2122}t a universally agreed-upon name for $%
\mathcal{V}_q $, it can be described using the following terminology:

\begin{itemize}
\item \textbf{Partial isometries with fixed range projection}: This
emphasizes that the final projector $q $ is fixed. 

\item \textbf{Equivalence class of partial isometries for $q$}: This
highlights the connection to Murray-von Neumann equivalence. 
\end{itemize}

\end{document}
